\begin{definition}
	Задание функций $F_j$ в $\R^{n+m},\ j=1,\ \dots,\ m$ означает задание отображения $\v F$, определённого на подмножестве $\R^{n+m}$ со значениями в $\R^n$. Непрерывность, дифференцируемость, непрерывная дифференцируемость этого отображения означает соответствующее свойство всех $F_j,\ j=1,\ \dots,\ m$.
\end{definition}
\begin{theorem}\nf{(Локальная обратимость отображений)}\label{27_th3}\\
	Если отображение $\v\phi$ непрерывно дифференцируемо в окрестности точки $\vx\n\in\R^n$, и якобиан $\frac{\mc D(\phi_1,\dots,\phi_n)}{\mc D(x_1,\dots,x_n)}(\vx\n)\ne0$, то в некоторой окрестности точки $\v y\n=\v\phi(\vx\n)$ определено обратное отображение $\v\phi^{-1}$, непрерывно дифференцируемое в этой окрестности.
\end{theorem}
\begin{proof}
	Положим $F_j(\vx,\v y):=\phi_j(\vx)-y_j,\ j=1,\ \dots,\ n$. Применим (\ref{26_th2}), поменяв ролями $\vx$ и $\v y$, тогда:
	\[(\fa\eps_1>0,\dots,\eps_n>0)(\ex\delta>0)\ K_{\delta,\v y\n}\ x_j=\psi_j(\v y)\lra F_j(\vx,\v y)=\phi_j(\vx)-y_j=0,\ j=1,\ \dots,\ n\]
	$\psi_j$ и является обратным отображением, получили требуемое.
\end{proof}
\begin{note}
	Чтобы найти частные производные неявных функций из (\ref{26_th2}) и компонент обратного отображения из (\ref{27_th3}), достаточно продифференцировать соответствующие тождества. Искомые частные производные находятся из получившихся систем линейных уравнений, определители которых совпадают с ненулевыми по условию якобианам.
\end{note}
\begin{note}
	Рассмотрим содержимое (\ref{27_th3}).
	\[\phi_j(\psi_1(y_1,\dots,y_n),\dots,\psi_n(y_1,\dots,y_n))=y_j,\ j=1,\ \dots,\ n\]
	\[\frac{\vd \phi_j}{\vd x_1}\cdot\frac{\vd \psi_1}{\vd y_k}+\dots+\frac{\vd \phi_j}{\vd k_n}\cdot\frac{\vd \psi_n}{\vd y_k}=\delta_{jk},\ j,\ k=1,\ \dots,\ n\]
	где $\delta_{jk}$ - символ Кронекера.\\
	То есть имеем:
	\begin{equation*}
		\begin{pmatrix}
			\frac{\vd \psi_1}{\vd y_1} & \dots & \frac{\vd \psi_1}{\vd y_n} \\
			\vdots & \ddots & \vdots \\
			\frac{\vd \psi_n}{\vd y_1} & \cdots & \frac{\vd \psi_n}{\vd y_n}
		\end{pmatrix}
		=
		\begin{pmatrix}
			\frac{\vd \phi_1}{\vd x_1} & \dots & \frac{\vd \phi_1}{\vd x_n} \\
			\vdots & \ddots & \vdots \\
			\frac{\vd \phi_n}{\vd x_1} & \cdots & \frac{\vd \phi_n}{\vd x_n}
		\end{pmatrix}
		^{-1}
	\end{equation*}
\end{note}
\begin{example}
	\begin{equation*}
		\begin{cases}
			x=r\cdot\cos\phi\\
			y=r\cdot\sin\phi
		\end{cases}
	\end{equation*}
	\[\frac{\mc D(x,y)}{\mc D(r,\phi)}=
	\begin{vmatrix}
		\cos\phi & -r\cdot\sin\phi \\
		\sin\phi & r\cdot\cos\phi
	\end{vmatrix}
	=r\ne0
	\]
	Заметим, что глобальной обратимости нет, так как образ точек с координатами $(r,\ \phi)$ и $(r,\ \phi+2\pi k)$ совпадает.
\end{example}
\subsection{Экстремумы функций многих переменных}
\begin{definition}
	Пусть $f$ определена в некоторой окрестности точки $\vx\in\R^n$. $\vx\n$ называется \textit{точкой локального максимума (минимума)} если:
	\[(\ex\delta>0)(\fa\vx\in U_\delta(\vx\n))\ f(\vx)\le f(\vx\n)\ (f(\vx)\ge f(\vx\n))\]
	В любом из этих случаем $\vx\n$ - точка локального экстремума функции $f$. Если в проколотой окрестности неравенства строгие, то точка $\vx\n$ - \textit{точка строго локального экстремума}.
\end{definition}
\begin{theorem}\nf{(Необходимое условие локального экстремума функций многих переменных)}\label{27_th1}\\
	Если $\vx\n$ является точкой локального экстремума дифференцируемой в этой точке функции $f(x)$, то:
	\[\frac{\vd f}{\vd x_j}(\vx\n)=0,\ j=1,\dots,n\]
\end{theorem}
\begin{proof}
	Зафиксируем все координаты $\vx\n$, кроме j-й. Обозначим полученную функцию $\psi_j(x_j):=f(x_1\n,\dots,x_{j-1}\n,x_j,x_{j+1}\n,\dots,x_n\n)$. \\
	Пусть $\vx\n$ - точка локального максимума. Тогда для $x_j\in(x_j\n-\delta,\ x_j\n+\delta)$:
	\[\psi_j(x_j)\le f(x_1\n,\dots,x_n\n)=\psi_j(x_j\n)\]
	Таким образом, $x_j\n$ - точка локального максимума функции $\psi_j$, дифференцируемой в точке $x_j\n$. Таким образом, $x_j\n$ - точка локального максимума функции $\psi_j$, дифференцируемой в $x_j\n$, а по теореме Ферма имеем:
	\[\frac{\vd f}{\vd x_j}(\vx\n)=\psi_j'(x_j\n)=0\]
	Получили требуемое.
\end{proof}
\begin{note}
	Заметим, что условие (\ref{27_th1}) эквивалентно записи $df(\vx\n)=0$.
\end{note}
\begin{example}
	Покажем, что равенства производных 0 недостаточно. Рассмотрим $f(x,y)=x^4-y^4$.\\
	\[\frac{\vd f}{\vd x}(0,0)=\frac{\vd f}{\vd y}(0,0)=0\]
	Если зафиксировать $y=0$, то точка $x=0$ - точка локального минимума, а если зафиксировать $x=0$, то $y=0$ - точка локального максимума.
\end{example}
\begin{definition}
	Если $df(\vx\n)\equiv0$, то $\vx\n$ называется \textit{стационарной точкой} функции f.
\end{definition}
\begin{note}
	\[d^2f(\vx\n)=\l( \frac\vd{\vd x_1}dx_1+\dots+\frac\vd{\vd x_n}dx_n \r)^2f(\vx\n)=\sum_{j,k=1}^n\frac{\vd^2f}{\vd x_j\cdot\vd x_k}(\vx\n)dx_jd_k\]
	Заметим, что для дважды дифференцируемое функции второй дифференциал является квадратичной формой.
\end{note}
\begin{theorem}\nf{(Достаточное условие локального экстремума)}\label{27_th2}\\
	Пусть f дважды дифференцируема в точке $\vx\n\in\R^n$ и дифференцируема в некоторой окрестности этой точки, причём $\vx\n$ - стационарная точка f. Тогда:
	\begin{enumerate}
		\item Если $d^2f(\vx\n)$ - положительно определённая квадратичная форма, то $\vx\n$ - точка строгого локального минимума f.
		\item Если $d^2f(\vx\n)$ - отрицательно определённая квадратичная форма, то $\vx\n$ - точка строгого локального максимума f.
		\item Если $d^2f(\vx\n)$ - знаконеопределённая квадратичная форма, то $\vx\n$ не является точкой локального экстремума.
	\end{enumerate}
\end{theorem}
\begin{proof}
	По формуле Тейлора с остаточным членом в форме Пеано:
	\[f(\vx)-f(\vx\n)=df(\vx\n)+\frac{d^2f(\vx\n)}{2}+\sigma(\rho^2),\ \rho\ra0\text{, где }\v{dx}=\vx-\vx\n,\ \rho=|\vx-\vx\n|\]
	\[d^2f(\vx\n)=\sum_{i,k=1}^n\frac{\vd^2f}{\vd x_i\vd x_j}(\vx\n)(x_i-x_i\n)(x_j-x_j\n)\]
	\[K(\v\xi)=\sum_{i,j=1}^n\frac{\vd^2f}{\vd x_i\vd x_j}(\vx\n)\xi_i\xi_j>0\text{ при }\v\xi>\v0\]
	K - квадратичная форма на компактном множестве. Рассмотрим её как функцию от $\xi$ на множестве S, где S:
	\[S^{n-1}=\{\v\xi\in\R^n:\xi_1^2+\dots+\xi_n^2\}\]
	$S^{n-1}$ - замкнуто и ограничено, значит ещё и компактно. Тогда по теореме Вейерштрасса:
	\[(\ex\v\xi_0\in S^{n-1})(\fa\v\xi\in S^{n-1})\ 0<m=K(\v\xi_0)\le K(\v\xi)\]
	Положим:
	\[x_i-x_i\n=\rho\cdot\frac{x_i-x_i\n}{\rho}=\rho\cdot\xi_i\]
	\[f(\vx)-f(\vx\n)=\rho^2\l(\frac{K(\v\xi)}2+\sigma(1)\r)\]
	Для достаточно малых $\rho:\ |\sigma(1)|\le\frac m2$, следовательно:
	\[(\fa \vx: 0<|\vx-\vx\n|<\delta)\ f(\vx)-f(\vx\n)\ge\frac{\rho^2m}4\]
	Получили, что точка $\vx\n$ является строгим локальным максимумом, первый пукнт доказан. Второй доказывается аналогично, теперь третий.
	\[\ex\v\xi\n[1],\v\xi\n[2],\ K(\v\xi\n[1])>0,\ K(\v\xi\n[2])<0,\ (\fa \rho>0)\ K(\rho\v\xi\n[1])>0,\ K(\rho\v\xi\n[2])<0\]
	Пусть $\vx\n[1]:=\rho\v\xi\n[1],\ \vx\n[2]:=\rho\v\xi\n[2]$, тогда:
	\[(\fa\delta>0)(\ex\vx\n[1],\vx\n[2],\ |\vx\n[i]-\vx\n|=\rho<\delta,\ i=1,2)\ K(\vx\n[i])=\rho^2K(\v\xi\n[i])>0,\ i=1,2,\text{ при }i=2\ <0\]
	\[f(\vx\n[1])-f(\vx\n)=\rho^2\l(\frac{K(\v\xi\n[1])}{2}+\sigma(1)\r)>0\]
	\[f(\vx\n[2])-f(\vx\n)=\rho^2\l(\frac{K(\v\xi\n[2])}{2}+\sigma(1)\r)<0\]
	Получили 2 сколь угодно близкие к $\vx\n$, значения функции в которых может быть как больше, так и меньше, значит $\vx\n$ не может быть точкой локального экстремума.\\
	Получили требуемое.
\end{proof}
\begin{corollary}
	Если $f(x,y)$ дважды дифференцируема в точке $(x_0,y_0)$ и дифференцируема в некоторой окрестности точки $(x_0,y_0)$, причём $(x_0,y_0)$ - стационарная точка $f$, то:
	\begin{enumerate}
		\item При $\l(f''_{xx}\cdot f''_{yy}-\l(f''_{xy}\r)^2\r)(x_0,y_0)>0$ точка $(x_0,y_0)$ является точкой строгого локального экстремума. Минимума, если $f''_{xx}(x_0,y_0)>0$ и максимума, если $f''_{xx}(x_0,y_0)<0$.
		\item При $\l(f''_{xx}\cdot f''_{yy}-\l(f''_{xy}\r)^2\r)(x_0,y_0)<0$ точка $(x_0,y_0)$ не является точкой локального экстремума функции f.
	\end{enumerate}
\end{corollary}
\begin{proof}
	Рассмотрим матрицу квадратичной формы для функции, определённой на $\R^2$:
	\begin{equation}\label{27_m3}
		\begin{vmatrix}
			f''_{xx} & f''_{xy} \\
			f''_{yx} & f''_{yy}
		\end{vmatrix}
	\end{equation}
	Воспользуемся критерием Сильвестра, получим требуемое.
\end{proof}
\begin{note}
	Матрицу (\ref{27_m3}) называют \textit{гессианом}.
\end{note}
\begin{definition}
	Пусть функции $f,\ \phi_1,\ \dots,\ \phi_m$ заданы на открытом множестве $G\subset\R^n,\ 1\le m<n$. $E=\{\vx\in G:\phi_i(\vx)=0,\ i=1,\ \dots,\ m\}$. Точка $\vx\n\in E$ называется \textit{точкой условного максимума (минимума)} функции $f$ при наличии условий связи (уравнений) $\phi_i(\vx)=0,\ i=1,\dots,\ m$, если:
	\[(\ex\delta>0)(\fa \vx\in U_\delta(\vx\n)\cap E)\ f(\vx)\le f(\vx\n)\ (f(\vx)\ge f(\vx\n))\]
	Заметим, что все эти условия автоматически выполняются, если точка $\vx\n$ является изолированной для множества $E$.
\end{definition}
\begin{note}
	\begin{equation}\label{27_m4}
		\begin{pmatrix}
			\frac{\vd \phi_1}{\vd x_1} & \dots & \frac{\vd \phi_1}{\vd x_n} \\
			\vdots & \ddots & \vdots \\
			\frac{\vd \phi_n}{\vd x_1} & \cdots & \frac{\vd \phi_n}{\vd x_n}
		\end{pmatrix}
	\end{equation}
	Как решать подобные вещи? Если функции $\phi_i$ непрерывно дифференцируемы в окрестности точки $\vx\n$ и ранг матрицы Якоби (\ref{27_m4}) в точке $\vx\n$ максимален, то можно выразить в окрестности точки $\vx\n$ $m$ переменных через оставшиеся.\\
	Будем считать, что $\ds\frac{\mc D(\phi_1,\dots,\phi_m)}{\mc D(x_1,\dots,x_m)}\ne0$ в точке $\vx\n$. Тогда система $\phi_i(\vx)=0,\ i=1,\dots,\ m\ \lra x_j=\psi_j(x_{m+1},\dots,x_n),\ j=1,\dots,m$.\\
	Обозначим $\t\vx\n:=(x_{m+1}\n,\dots,x_n\n)$.\\
	Возможность почленного дифференцирования равенств $x_i=\phi_i(\vx),\ i=1,\dots,\ m$ и $x_j=\psi_j(\t\vx\n),\ j=1,\dots,\ m$ упоминалось в замечании после (\ref{26_th2}). Для этого надо в ходе доказательства теоремы 2 по индукции проверить эту возможность.\\
	Была система:
	\[
	\begin{cases*}
		F_1(x_1,\dots,x_n,y_1,\dots,y_m)=0\\
		\dots\\
		F_m(x_1,\dots,x_n,y_1,\dots,y_m)=0
	\end{cases*}
	\lra
	\begin{cases}
		y_1=\phi_1(x_1,\dots,x_n)\\
		\dots\\
		y_m=\phi_m(x_1,\dots,x_n)
	\end{cases}
	\]
	Хотим доказать, что:
	\[
	\begin{cases}
		dF_1(x_1,\dots,x_n,y_1,\dots,y_m)=0\\
		\dots\\
		dF_m(x_1,\dots,x_n,y_1,\dots,y_m)=0
	\end{cases}
	\lra
	\begin{cases}
		dy_1=d\phi_1(x_1,\dots,x_n)\\
		\dots\\
		dy_m=d\phi_m(x_1,\dots,x_n)
	\end{cases}
	\]
	По гипотезе индукции есть системы:
	\[
		\begin{cases*}
			F_1(x_1,\dots,x_n,y_1,\dots,y_m)=0\\
			\dots\\
			F_{m-1}(x_1,\dots,x_n,y_1,\dots,y_m)=0
		\end{cases*}
		\lra
		\begin{cases}
			y_1=\phi_1(x_1,\dots,x_n)\\
			\dots\\
			y_{m-1}=\phi_m(x_1,\dots,x_n)
		\end{cases}
	\]
	\[
	\begin{cases}
		dF_1(x_1,\dots,x_n,y_1,\dots,y_m)=0\\
		\dots\\
		dF_{m-1}(x_1,\dots,x_n,y_1,\dots,y_m)=0
	\end{cases}
	\lra
	\begin{cases}
		dy_1=d\phi_1(x_1,\dots,x_n)\\
		\dots\\
		dy_{m-1}=d\phi_m(x_1,\dots,x_n)
	\end{cases}
	\]
	В доказательстве в $F_m(x_1,\dots,x_n,y_1,\dots,y_m)$ подставляли $\Phi_1,\dots,\Phi_{m-1}$ и получали:
	\[F_m(x_1,\dots,x_n,\Phi_1(x_1,\dots,x_n,y_m),\dots,\ \Phi_{m-1}(x_1,\dots,x_n,y_m),y_m)=0\lra y_m=\phi_m(x_1,\dots,x_n)\]
	В силу инвариантности формы первого дифференциала:
	\[dF_m(x_1,\dots,x_n,y_1,\dots,y_m)=0=dF_m(\dots)\lra dy_m=d\phi_m(x_1,\dots,x_n)\]
	\[\phi_i(x_1,\dots,x_n):=\Phi_i(x_1,\dots,x_n,\phi_m(x_1,\dots,x_n))\]
	\[dy_i=d\Phi_i(x_1,\dots,x_n,y_m)=d\phi_i(x_1,\dots,x_n)\]
	Получили требуемое.\\
	Вернёмся к задаче. Если 2 способа:
	\begin{enumerate}
		\item Найти в явном виде решение системы $\phi_i(\vx)=0,\ i=1,\dots,\ m$, подставить их в f, то есть рассмотреть функцию $F(x_{m+1},\dots,x_n)=f(\psi_1(\t\vx),\dots,\psi_m(\t\vx),\t\vx)$, где $\t\vx=(x_{m+1},\dots,x_n)$. Тогда задача на условный экстремум равносильна задаче на локальный (безусловный) экстремум для $F$.
		\item Рассмотреть функцию Лагранжа $L(x_1,\dots,x_n,\la_1,\dots,\la_m):=f(\vx)+\sum_{i=1}^m\la_i\phi_i(\vx)$, где $\la_i,\ i=1,\dots,m$ - множители Лагранжа.
	\end{enumerate}
\end{note}
\begin{theorem}\nf{(Необходимое условие условного экстремума)}\label{27_th5}\\
	Если дифференцируемая в точке $\vx\n$ функция f имеет в точке $\vx\n$ условный экстремум при наличии уравнений связей $\phi_i(\vx)=0,\ i=1,\dots,m$, заданными непрерывно дифференцируемыми в окрестности точки $\vx\n$ функциями $\phi_i$ такими, что ранг матрицы Якоби (\ref{27_m4}) равен m ($1\le m<n$), то существует единственный набор $(\la_1\n,\dots,\la_m\n)$ такой, что $(\vx\n,\v\la)$ является стационарной точкой функции Лагранжа $L(\vx,\v\la)=f(\vx)+\sum_{i=1}^n\la_i\phi_i(\vx)$.
\end{theorem}
\begin{proof}
	Задача на условный экстремум равносильна задаче на безусловный экстремум функции $F(\t\vx)=f(\psi_1(\t\vx),\dots,\psi_m(\tvx),\tvx)$, где функции $\psi_i$ определяются из равенств $x_j=\psi_j(x_{m+1},\dots,x_n),\ j=1,\dots,m$, которые равносильны $\phi_i(\vx)=0,\ i=1,\dots,m$ в окрестности точки $\vx\n$.\\
	Значит, по (\ref{27_th1}), $\tvx\n$ является стационарной точкой функции $F$, то есть $dF(\tvx\n)=0$.\\
	Посчитаем этот дифференциал. Из инвариантности формы первого дифференциала:
	\begin{equation}\label{28_eq2}
		\sum_{i=1}^n\frac{\vd f}{\vd x_i}(\vx\n)dx_i=0
	\end{equation}
	где $dx_1,\dots,dx_m$ - дифференциалы функций $\psi_1,\dots,\psi_m$.\\
	Воспользуемся почленным дифференцированием из замечания выше:
	\[dx_j=d\psi_j(\tvx\n),\ j=1,\dots,m\lra d\phi_j(\vx\n)=0,\ j=1,\dots,m\]
	\[d\phi_j(\vx\n)=\sum_{i=1}^n\frac{\vd \phi_j}{\vd x_i}dx_i\]
	Домножим на $\la_j$ и сложим с $\ds\sum_{i=1}^n\frac{\vd f}{vd x_i}(\vx\n)dx_i=0$:
	\[0=\sum_{i=1}^n\l(\frac{\vd f}{\vd x_i}(\vx\n)+\sum_{j=1}^m\la_j\frac{\vd \phi_j}{\vd x_j}(\vx\n)\r)dx_i=\sum_{i=1}^n\frac{\vd L}{\vd x_i}(\vx\n,\v\la)dx_i\]
	Найдём такие $\la_j$, что $\ds\frac{\vd L}{\vd x_i}(\vx\n,\v\la)=0,\ i=1,\dots,m$. Это можно сделать, так как определитель соответствующей системы линейных уравнений (\ref{27_m4}) не равен 0. Обозначим решение этой системы за $\v\la\n$, тогда:
	\[\sum_{i=1}^n\frac{\vd L}{\vd x_i}(\vx\n,\v\la)dx_i=\sum_{i=m+1}^n\frac{\vd L}{\vd x_i}(\vx\n,\v\la\n)dx_i=0\]
	Значит, все частные производные функции Лагранжа по всем переменным равны 0. \\
	Получили требуемое.
\end{proof}